\chapter{Buenas prácticas de seguridad}
\noindent \textbf{Recomendación 1:} Al utilizar \textit{wearables} (pulseras o relojes inteligentes), se debe configurar los mecanismos de protección de la propia aplicación que gestiona el dispositivo y revisar las opciones de privacidad de la red social con la que se sincroniza, ya que estos sensores capturan datos biométricos altamente sensibles sobre la salud. \\
\textbf{Ficha número:} 18

\vspace{1cm}

\noindent \textbf{Recomendación 2:} En servicios de mensajería instantánea, se recomienda hacer uso de la opción de ``chat privado'' o ``secreto'' para evitar que personas ajenas a la conversación puedan espiarla, además de asegurar que el intercambio de mensajes esté siempre cifrado. \\
\textbf{Ficha número:} 10

\vspace{1cm}

\noindent \textbf{Recomendación 3:} Al utilizar navegadores, se debe instalar un ``verificador de páginas web'' (generalmente proporcionado por los principales antivirus) para obtener una capa adicional de seguridad que valide la reputación de los sitios antes de interactuar con ellos. \\
\textbf{Ficha número:} 8